\cotefile{157-1}
\scanpage{114}
\mspage{120}
Par \unsure{contre}, on a

\uline{Prop. 8} Soit $f:X\to Y$ dans \unsure{Diag}. Conditions équivalentes

a) $f$ est une $\unsure{\mathbb D|D}$-fibration

b) \unsure{Pour} tout $p:Y''\unsure{\mathord{-}|\to}Y'$ sur $Y$, si $p$ est une $\mathbb D$-équiv. $1$-complète \unsure{(}déf. 1 p.~113), alors $q:X''-X'$ est une $\mathbb D$-équivalence.

% diagram in the left margin beside b)
\diagram{\begin{tikzcd}[ampersand replacement=\&,column sep=small,row sep=small]
X\arrow[d,"f"'] \& X'\arrow[l]\arrow[d,"f'"'] \& X''\arrow[l,"q"']\arrow[d,"f''"]\\
Y \& Y'\arrow[l] \& Y''\arrow[l,"p"]
\end{tikzcd}}

\uline{Cor.} Supposons que dans Diag, toute $\mathbb D$-équivalence soit $1$-complète. (NB Je ne connais que le dérivateur HOT pour satisfaire à cette condition, qui \unsure{paraît} quasiment \unsure{caractéristique} de HOT.) Alors $f:X-Y$ est une $\mathbb D$-\deleted{équiv.}fibration \unsure{ssi} pour tout $p:Y''-Y'$ sur $Y$, si $p$ est une $\mathbb D$-équiv., alors $q:X''-X'$ aussi. Dans ce cas ``exceptionnel'' parmi les dérivateurs, la notion de $\mathbb D$-fibration est stable par composition (ce qui est \unsure{tautologique|tout le gag} sur le critère que je viens de donner.)

On va renforcer la notion de $\mathbb D$-fibration, pour la rendre plus souple, et notamment stable par composition.

\uline{Proposition 9} Soit $f:X-Y$ dans Diag. Conditions équivalentes :

a) $\forall\,y\in Y,\quad X_y-X_{/y}$ est une $\mathbb D$-\uline{équiv. \unsure{cohomologique}} (déf. 1 p.~114)\unsure{''|,} et \unsure{de|en} $=$ pour $X_{/y}-X_{/y'}$ pour toute flèche $y\longrightarrow y'$ dans $Y$. Et ces conditions sur $f$ restent valables après tout chgt de base (\uline{NB} il suffit que la première de ces conditions le reste.)\unsure{.}
% archive sheet number: 114
